\documentclass[12pt,a4paper]{report}
\usepackage[utf8]{inputenc}
\usepackage[T1]{fontenc}
\usepackage[french]{babel}
\usepackage{amsmath, amssymb, amsthm, amsfonts}
\usepackage{geometry}
\usepackage{xcolor}
\usepackage{listings}
\usepackage{tikz}
\geometry{hmargin=2.5cm,vmargin=2.5cm}

% Environnements de théorèmes (style Ch. 3)
\newtheorem{theoreme}{Théorème}[section]
\newtheorem{proposition}[theoreme]{Proposition}
\newtheorem{lemme}[theoreme]{Lemme}
\newtheorem{corollaire}[theoreme]{Corollaire}
\theoremstyle{definition}
\newtheorem{definition}{Définition}[section]
\newtheorem{exemple}{Exemple}[section]
\theoremstyle{remark}
\newtheorem{remarque}{Remarque}[section]
\newtheorem{propriete}{Propriété}[section]

% Code Python
\lstset{
    language=Python,
    basicstyle=\ttfamily\small,
    keywordstyle=\color{blue}\bfseries,
    stringstyle=\color{red},
    commentstyle=\color{green!60!black}\itshape,
    numbers=left,
    numberstyle=\tiny\color{gray},
    stepnumber=1,
    frame=single,
    breaklines=true
}

\title{Algèbre 1 : Chapitre IV --- Structures algébriques}
\author{Aymen Ben Brik}
\date{}

\begin{document}

\maketitle

\setcounter{chapter}{3}
\chapter{Structures algébriques}

Ce chapitre introduit les structures fondamentales de l'algèbre : groupes, anneaux et corps. Ces structures unifient sous un même formalisme l'arithmétique entière, le calcul matriciel, les permutations, l'arithmétique modulaire de la cryptographie, ainsi que de nombreuses transformations en informatique graphique et en science des données. La force de cette abstraction est qu'un théorème démontré sur une structure s'applique automatiquement à toutes ses incarnations concrètes.

\section{Structure de groupe}

\subsection{Motivation}
Un \emph{groupe} est l'abstraction la plus simple qui capture la notion de symétrie réversible. Que vous mélangiez un paquet de cartes, fassiez tourner un cube de Rubik, chiffriez un message par un protocole asymétrique (Diffie-Hellman, RSA), ou appliquiez une transformation géométrique en infographie, vous manipulez sans le savoir un groupe. Cette unification est précieuse : un théorème démontré une fois sur les groupes (par exemple « tout sous-groupe d'un groupe cyclique est cyclique ») s'applique automatiquement à des dizaines de situations très différentes.

\subsection{Approche par problème : l'horloge et le chiffrement de César}
\textbf{Problème :} Dans le chiffrement de César de pas $k$, on remplace chaque lettre par celle qui se trouve $k$ positions plus loin dans l'alphabet (modulo 26).
\begin{enumerate}
    \item Si je chiffre avec $k_1 = 5$ puis $k_2 = 9$, par quel décalage déchiffrer ?
    \item Si je chiffre par $k = 13$, que se passe-t-il en chiffrant deux fois ?
    \item Quelles propriétés possède l'opération « composer deux décalages » sur $\{0, 1, \dots, 25\}$ ?
\end{enumerate}

\textbf{Résolution :}
La composition correspond à l'addition modulo $26$. Donc $k_1 + k_2 = 14$ : pour déchiffrer il faut décaler de $-14 \equiv 12 \pmod{26}$.
Pour $k = 13$ : $13 + 13 = 26 \equiv 0$. Chiffrer deux fois revient à l'identité, c'est ROT13.
La loi $+ \pmod{26}$ est associative, possède un neutre ($0$) et tout élément $k$ a un inverse $26 - k$. C'est exactement la structure de \emph{groupe} sur l'ensemble $\mathbb{Z}/26\mathbb{Z}$.

\subsection{Loi de composition interne et définition du groupe}

\begin{definition}[Loi de composition interne]
Une \textbf{loi de composition interne} (l.c.i) sur un ensemble non vide $G$ est une application $\star : G \times G \to G$ qui à un couple $(a, b)$ associe l'élément noté $a \star b$.
\end{definition}

\begin{definition}[Groupe]
Un couple $(G, \star)$ est un \textbf{groupe} si $\star$ est une l.c.i sur $G$ vérifiant :
\begin{enumerate}
    \item \textbf{(Associativité)} $\forall a, b, c \in G,\ (a \star b) \star c = a \star (b \star c)$ ;
    \item \textbf{(Élément neutre)} $\exists e \in G,\ \forall a \in G,\ a \star e = e \star a = a$ ;
    \item \textbf{(Inverse)} $\forall a \in G,\ \exists a' \in G,\ a \star a' = a' \star a = e$.
\end{enumerate}
Si de plus $\forall a, b \in G,\ a \star b = b \star a$, le groupe est dit \textbf{abélien} (ou commutatif).
\end{definition}

\begin{proposition}[Unicité du neutre et de l'inverse]
Dans tout groupe $(G, \star)$ :
\begin{itemize}
    \item l'élément neutre $e$ est unique ;
    \item pour tout $a \in G$, l'inverse $a'$ est unique. On le note $a^{-1}$ (notation multiplicative) ou $-a$ (notation additive).
\end{itemize}
\end{proposition}

\begin{proof}
Si $e$ et $e'$ sont neutres : $e = e \star e' = e'$. Si $a'$ et $a''$ sont inverses de $a$ : $a' = a' \star e = a' \star (a \star a'') = (a' \star a) \star a'' = e \star a'' = a''$.
\end{proof}

\begin{exemple}[Groupes usuels]
\begin{itemize}
    \item $(\mathbb{Z}, +)$, $(\mathbb{Q}, +)$, $(\mathbb{R}, +)$, $(\mathbb{C}, +)$ sont des groupes abéliens.
    \item $(\mathbb{Q}^*, \times)$, $(\mathbb{R}^*, \times)$, $(\mathbb{C}^*, \times)$ sont des groupes abéliens.
    \item $(\mathbb{N}, +)$ \emph{n'est pas} un groupe : pas d'inverse pour $n \ge 1$.
    \item L'ensemble $S_n$ des permutations de $\{1, \dots, n\}$ muni de la composition est un groupe, non abélien dès $n \ge 3$.
    \item $GL_n(\mathbb{R})$, l'ensemble des matrices réelles $n \times n$ inversibles muni du produit, est un groupe (non abélien dès $n \ge 2$).
\end{itemize}
\end{exemple}

\subsection{Sous-groupes}

\begin{definition}[Sous-groupe]
Soit $(G, \star)$ un groupe d'élément neutre $e$. Une partie $H \subseteq G$ est un \textbf{sous-groupe} de $G$, noté $H \le G$, si :
\begin{enumerate}
    \item $e \in H$ ;
    \item $\forall a, b \in H,\ a \star b \in H$ \quad (stabilité par la loi) ;
    \item $\forall a \in H,\ a^{-1} \in H$ \quad (stabilité par inverse).
\end{enumerate}
\end{definition}

\begin{proposition}[Caractérisation pratique d'un sous-groupe]
$H \subseteq G$ est un sous-groupe de $G$ si et seulement si :
\[ H \neq \emptyset \quad \text{et} \quad \forall a, b \in H,\ a \star b^{-1} \in H. \]
\end{proposition}

\begin{proof}
($\Rightarrow$) immédiat. ($\Leftarrow$) $H$ non vide contient un élément $a$, donc $a \star a^{-1} = e \in H$. Puis $\forall b \in H$, $e \star b^{-1} = b^{-1} \in H$ donc $H$ est stable par inverse. Enfin $\forall a, b \in H$, $a \star (b^{-1})^{-1} = a \star b \in H$ donc $H$ est stable par la loi.
\end{proof}

\begin{remarque}
Tout sous-groupe contient le neutre et est lui-même un groupe pour la loi induite. Pour tout groupe $G$, $\{e\}$ et $G$ sont les sous-groupes \emph{triviaux}.
\end{remarque}

\subsection{Les sous-groupes de $(\mathbb{Z}, +)$}

Les sous-groupes du groupe additif $\mathbb{Z}$ se décrivent intégralement.

\begin{theoreme}[Sous-groupes de $\mathbb{Z}$]
Soit $H$ un sous-groupe de $(\mathbb{Z}, +)$. Il existe un unique entier $n \in \mathbb{N}$ tel que
\[ H = n\mathbb{Z} = \{nk : k \in \mathbb{Z}\}. \]
\end{theoreme}

\begin{proof}
Si $H = \{0\}$, on prend $n = 0$. Sinon $H$ contient un entier non nul, donc un entier strictement positif (en prenant son opposé si besoin). Posons $n = \min(H \cap \mathbb{N}^*)$.

\emph{Inclusion $n\mathbb{Z} \subseteq H$ :} par stabilité par addition et passage à l'opposé, $n \in H$ entraîne $kn \in H$ pour tout $k \in \mathbb{Z}$.

\emph{Inclusion $H \subseteq n\mathbb{Z}$ :} soit $h \in H$. La division euclidienne donne $h = nq + r$ avec $0 \le r < n$. Comme $h \in H$ et $nq \in H$, on a $r = h - nq \in H$. Mais $r < n$ et $n$ est le plus petit élément strictement positif de $H$, donc $r = 0$. Ainsi $h = nq \in n\mathbb{Z}$.

L'unicité est claire : $n\mathbb{Z} = m\mathbb{Z}$ avec $n, m \ge 0$ implique $n = m$ (plus petit élément $> 0$ commun).
\end{proof}

\begin{remarque}
Ce résultat est central. Il dit que tous les sous-groupes de $\mathbb{Z}$ sont \emph{cycliques} (engendrés par un seul élément, voir ci-dessous). On retrouvera ce phénomène dans tout groupe cyclique.
\end{remarque}

\subsection{Sous-groupe engendré, groupes monogènes et cycliques}

\begin{definition}[Sous-groupe engendré par un élément]
Soit $(G, \star)$ un groupe et $a \in G$. On note $\langle a \rangle$ le \textbf{sous-groupe engendré par $a$} :
\[ \langle a \rangle = \{a^n : n \in \mathbb{Z}\} \]
en notation multiplicative, où $a^0 = e$, $a^n = \underbrace{a \star \cdots \star a}_{n \text{ fois}}$ et $a^{-n} = (a^{-1})^n$ pour $n \ge 1$. En notation additive, $\langle a \rangle = \{na : n \in \mathbb{Z}\}$. C'est le plus petit sous-groupe de $G$ contenant $a$.
\end{definition}

\begin{definition}[Groupe monogène, groupe cyclique]
Un groupe $G$ est \textbf{monogène} s'il existe $a \in G$ tel que $G = \langle a \rangle$. Un tel $a$ est appelé \emph{générateur}. Si de plus $G$ est fini, on dit que $G$ est \textbf{cyclique}.
\end{definition}

\begin{exemple}[Premiers exemples]
\begin{itemize}
    \item $(\mathbb{Z}, +)$ est monogène, engendré par $1$ (ou par $-1$). Il est infini, donc non cyclique.
    \item $(\mathbb{Z}/n\mathbb{Z}, +)$ est cyclique d'ordre $n$, engendré par $\bar{1}$. Plus généralement, $\bar{k}$ est générateur si et seulement si $\mathrm{pgcd}(k, n) = 1$.
    \item Le groupe $S_3$ a $6$ éléments mais aucun n'engendre les six : il n'est pas cyclique.
\end{itemize}
\end{exemple}

\subsection{Ordre d'un élément}

\begin{definition}[Ordre d'un élément]
Soit $(G, \star)$ un groupe (noté multiplicativement) et $a \in G$. L'\textbf{ordre} de $a$, noté $o(a)$, est :
\begin{itemize}
    \item $o(a) = +\infty$ si $\langle a \rangle$ est infini ;
    \item $o(a) = $ le plus petit entier $n \ge 1$ tel que $a^n = e$, sinon.
\end{itemize}
Lorsque $o(a)$ est fini, on a $|\langle a \rangle| = o(a)$.
\end{definition}

\begin{proposition}[Caractérisation de l'ordre]
Soit $a \in G$ d'ordre fini $n$. Pour tout $k \in \mathbb{Z}$ :
\[ a^k = e \iff n \mid k. \]
\end{proposition}

\begin{proof}
Division euclidienne $k = nq + r$ avec $0 \le r < n$ : $a^k = (a^n)^q \star a^r = e^q \star a^r = a^r$. Donc $a^k = e \iff a^r = e$. Comme $0 \le r < n$ et $n$ est le plus petit entier $\ge 1$ vérifiant $a^n = e$, ceci équivaut à $r = 0$, soit $n \mid k$.
\end{proof}

\begin{exemple}[Ordres dans $(\mathbb{Z}/12\mathbb{Z}, +)$]
On cherche le plus petit $k \ge 1$ tel que $k \cdot a \equiv 0 \pmod{12}$, ce qui donne $k = 12 / \mathrm{pgcd}(a, 12)$.
\begin{center}
\begin{tabular}{c|cccccccccccc}
$a$ & 0 & 1 & 2 & 3 & 4 & 5 & 6 & 7 & 8 & 9 & 10 & 11 \\
\hline
$o(a)$ & 1 & 12 & 6 & 4 & 3 & 12 & 2 & 12 & 3 & 4 & 6 & 12 \\
\end{tabular}
\end{center}
On observe que $o(a) \mid 12$ pour tout $a$ : c'est un cas particulier du théorème de Lagrange (à venir au §1.b).
\end{exemple}

\subsection{Exemples résolus (Difficulté ascendante)}

\textbf{Exemple 1 (Vérification d'un groupe --- Facile).}
Montrer que $(\mathbb{R}^* , \times)$ est un groupe abélien.

\textit{Correction :} la multiplication réelle est interne sur $\mathbb{R}^*$ (un produit de réels non nuls est non nul), associative et commutative. L'élément $1$ est neutre. Tout $x \in \mathbb{R}^*$ admet l'inverse $1/x \in \mathbb{R}^*$.

\medskip

\textbf{Exemple 2 (Sous-groupe par caractérisation --- Intermédiaire).}
Montrer que $H = \{2^n : n \in \mathbb{Z}\}$ est un sous-groupe de $(\mathbb{R}^*_+, \times)$.

\textit{Correction :} $H \neq \emptyset$ (il contient $1 = 2^0$). Pour $a = 2^n$ et $b = 2^m$ dans $H$, on a $a \times b^{-1} = 2^n \times 2^{-m} = 2^{n-m} \in H$ puisque $n - m \in \mathbb{Z}$. La caractérisation pratique conclut : $H \le \mathbb{R}^*_+$.

\medskip

\textbf{Exemple 3 (Calcul d'ordre dans un groupe modulaire --- Difficile).}
Calculer l'ordre de $\bar{3}$ dans $((\mathbb{Z}/14\mathbb{Z})^*, \times)$.

\textit{Correction :} on calcule les puissances modulo $14$ :
$3^1 = 3,\ 3^2 = 9,\ 3^3 = 27 \equiv 13 \equiv -1,\ 3^4 \equiv -3 \equiv 11,\ 3^5 \equiv 33 \equiv 5,\ 3^6 \equiv 15 \equiv 1 \pmod{14}$.
Donc $o(\bar{3}) = 6$. Le sous-groupe engendré est $\langle \bar{3} \rangle = \{1, 3, 9, 13, 11, 5\}$, qui contient les six éléments inversibles modulo $14$. Donc $\bar{3}$ est un générateur de $(\mathbb{Z}/14\mathbb{Z})^*$ (qui est cyclique d'ordre $\varphi(14) = 6$).

\subsection{Exercices pratiques avec Python}
On peut tester sur machine la structure de groupe de $\mathbb{Z}/n\mathbb{Z}$ : calculer l'ordre d'un élément, énumérer le sous-groupe engendré, lister tous les sous-groupes.

\begin{lstlisting}
from math import gcd

def ordre_additif(a, n):
    """Ordre de a dans (Z/nZ, +) : plus petit k>=1 tel que k*a == 0 mod n."""
    return n // gcd(a, n)

def sous_groupe_engendre_additif(a, n):
    """Sous-groupe engendre par a dans (Z/nZ, +)."""
    H = set()
    x = 0
    while x not in H:
        H.add(x)
        x = (x + a) % n
    return sorted(H)

def tous_les_sous_groupes_de_ZnZ(n):
    """Les sous-groupes de Z/nZ sont en bijection avec les diviseurs de n."""
    return {d: sous_groupe_engendre_additif(d, n) for d in range(1, n+1) if n % d == 0}

# Test sur Z/12Z
print("Ordres dans Z/12Z :")
for a in range(12):
    print(f"  o({a}) = {ordre_additif(a, 12) if a != 0 else 1}")

print("\nSous-groupes de Z/12Z (un par diviseur de 12) :")
for d, H in tous_les_sous_groupes_de_ZnZ(12).items():
    print(f"  <{d}> = {H}  (|H| = {len(H)})")
\end{lstlisting}

\subsection{Exercices à faire}
\begin{enumerate}
    \item \textbf{Exercice 1.} Soit $G = \{1, -1, i, -i\} \subset \mathbb{C}^*$. Montrer que $(G, \times)$ est un groupe cyclique. Donner ses générateurs et tracer sa table.
    \item \textbf{Exercice 2.} Soit $(G, \star)$ un groupe et $H_1, H_2$ deux sous-groupes de $G$. Montrer que $H_1 \cap H_2$ est un sous-groupe. La réunion $H_1 \cup H_2$ est-elle toujours un sous-groupe ? (Donner un contre-exemple.)
    \item \textbf{Exercice 3.} Lister tous les sous-groupes de $(\mathbb{Z}/18\mathbb{Z}, +)$ en justifiant. Combien y en a-t-il ?
    \item \textbf{Exercice 4.} Calculer l'ordre de $\bar{2}$, $\bar{3}$ et $\bar{5}$ dans $((\mathbb{Z}/11\mathbb{Z})^*, \times)$. Lequel est un générateur ?
    \item \textbf{Exercice 5 (synthèse).} Montrer que dans tout groupe $G$, si $a \in G$ vérifie $a^2 = e$, alors $\langle a \rangle = \{e, a\}$ et $o(a) \in \{1, 2\}$.
\end{enumerate}

\paragraph{Suite (§1.b).} Théorème de Lagrange, morphismes de groupes, le groupe symétrique $S_n$ et le groupe $\mathbb{Z}/n\mathbb{Z}$ (étude détaillée).

\end{document}
